% BT1202 -- R3 continuum hypothesis checklist insert
% Generated from the page-36+ photonic-holonet consolidation pass.

\subsection{R3 as a finite-to-continuum hypothesis checklist}
\label{sec:bt1202-r3-checklist}

The third physical residual should be read as a theorem-governed limit problem,
not as a new architectural gap.  Appendix~B separates two continua.  The
holonet's own limit is symplectic: finite $W(3,q)$ data pass to metaplectic and
continuous-variable oscillator computation.  The spacetime limit is metric: the
K3/spectral-action side must converge to a smooth Riemannian/Lorentzian field
geometry.  The architecture is complete before taking either limit; R3 asks
whether the metric continuum hypotheses hold for the chosen refinement family.

For reviewer-facing use, R3 is the following checklist.

\begin{center}
\small
\begin{tabular}{p{0.27\linewidth}p{0.43\linewidth}p{0.21\linewidth}}
\toprule
Hypothesis & What must be verified & Failure mode \\
\midrule
Shape-regular refinement & The finite K3/spacetime approximants refine without
collapsing simplex quality, injectivity scale, or incidence valence. & No stable
metric continuum. \\
Bounded curvature energy & Discrete curvature, heat-kernel coefficients, and
spectral-action moments remain uniformly bounded under refinement. & Curvature
concentrates into lattice artifacts. \\
Spectral moment matching & The low-order spectral-action coefficients converge
to the intended Einstein--Yang--Mills terms, with the known $E_8$/K3 gauge
lattice fixed. & Correct integers but wrong continuum action. \\
Gauge-field convergence & Holonomy/curvature data converge as connections, not
only as pointwise group labels. & Gate holonomy survives while spacetime gauge
field does not. \\
Operator convergence & Dirac/Laplacian approximants converge in an appropriate
strong-resolvent, spectral, or quantum-Gromov sense. & Spectrum fits finite
checks but has no continuum operator. \\
Metric-propinquity control & The finite metric spaces form a Cauchy family in a
metric strong enough to preserve algebraic and spectral structure. & The limit is
nonunique or topology-dependent. \\
Scale separation & The symplectic computer continuum and the metric spacetime
continuum remain coupled but not conflated. & A polar-graph limit is mistakenly
claimed to be spacetime. \\
\bottomrule
\end{tabular}
\end{center}

\begin{principle}[R3 is a convergence residual]
\label{pr:bt1202-r3-convergence-residual}
R3 is not the problem of finding more holonet architecture.  It is the problem
of proving that the already-fixed finite geometry admits a metric continuum
limit satisfying the listed compactness, spectral, gauge, and scale-separation
hypotheses.  The symplectic side supplies the computer; the metric side must
supply spacetime.
\end{principle}

This is the clean boundary.  Failure of any checklist row would falsify the
spacetime-continuum reading, but it would not erase the finite holonet
architecture: $W(3,3)$ would still define the symplectic quantum computer.  A
successful R3 proof would instead show that the same finite seed has two
compatible limits, a metaplectic/CV computation limit and a K3/spectral-action
metric limit.
